\documentclass[12pt, letterpaper]{article}
\usepackage[margin=1in]{geometry}
\usepackage{enumitem}
\usepackage{parskip}
\usepackage{hyperref}
\usepackage{titlesec}
\usepackage{booktabs}

\hypersetup{
  colorlinks=true,
  linkcolor=black,
  urlcolor=black
}

\titleformat{\section}{\large\bfseries}{}{0em}{}[\titlerule]

\title{\textbf{International Student AI Influencer}\\[0.4em]
\large Project Report --- Gen AI \& Social Media}
\author{
  Joaquin Sardon \and
  Kulbir Kaur \and
  Melissa Lee \and
  Shivani Kabra \and
  Gerald Velasquez
}
\date{May 2026}

\begin{document}
\maketitle
\thispagestyle{empty}
\newpage

% ─────────────────────────────────────────────────────────────────────────────
\section{Executive Summary}
% ─────────────────────────────────────────────────────────────────────────────

The International Student AI Influencer is an AI-powered multi-agent pipeline that
transforms scattered civic and financial information into platform-ready social media
content for international students. Built for the NYC Build With AI Hackathon 2026,
the system addresses the fragmented information landscape facing international students
navigating housing, taxes, banking, immigration, and city services. The pipeline
automates the full content lifecycle---from source retrieval to post generation, image
creation, video production, and quality control---while surfacing important risks around
trust, accuracy, and sensitive topic handling. Although prototyped for New York City
using NYC Open Data, the architecture is designed as a reusable pipeline adaptable to
other cities or university contexts by swapping data connectors.

% ─────────────────────────────────────────────────────────────────────────────
\section{Problem \& Motivation}
% ─────────────────────────────────────────────────────────────────────────────

International students arriving in the United States face a fragmented and often
overwhelming information landscape. Civic resources covering housing, employment, taxes,
immigration, banking, safety, and city services are dispersed across government portals,
institutional websites, and informal peer networks. The slide deck identifies this
fragmentation as the core problem: no single, trustworthy, student-facing channel
consolidates this information in an accessible, platform-native format.

The repository reinforces this motivation technically. Agent 1 aggregates from
heterogeneous sources including live Google Search results (with citations), three NYC
Open Data datasets (job postings, CUNY enrollment figures, and living wage benchmarks),
eight financial RSS feeds spanning the IRS, CFPB, NASFAA, NYT Business, SoFi, Prodigy,
Earnest, and MPOWER, and a curated set of seven F-1/J-1 visa tax scenarios. This breadth
of data sources directly mirrors the breadth of the problem: students require one
coherent, trustworthy voice across many interconnected domains.

The slide deck and repository are in strong agreement here. The slide deck frames the
problem from a student experience perspective; the repository demonstrates that the
technical data model was designed to match that scope from the ground up.

% ─────────────────────────────────────────────────────────────────────────────
\section{Technical Architecture}
% ─────────────────────────────────────────────────────────────────────────────

The system is built on a five-agent sequential pipeline with an optional sixth agent
for financial scenario resolution. The interface layer is a React 19 + Vite +
TypeScript single-page application backed by a FastAPI REST server. All generative
models are provided by Google Vertex AI.

\textbf{Agent Pipeline:}
\begin{itemize}[leftmargin=1.5em]
  \item \textbf{Agent 1 --- Source Retrieval:} Queries Google Search Grounding, NYC
    Open Data, financial RSS feeds, and a static tax knowledge seed. Documents are
    chunked (400 characters, 50-character overlap), embedded with
    \texttt{gemini-embedding-001}, and ranked via in-memory cosine similarity with
    topic weight boosting.
  \item \textbf{Agent 2 --- Content Generator:} Detects trending angles, generates
    post text, hashtags, and calls-to-action, and produces a five-scene video storyboard
    with an SSML voice script using Gemini 2.5 Flash.
  \item \textbf{Creative Storyteller (Output Layer):} Generates social media images
    via an Imagen fallback chain (4.0-ultra $\to$ 4.0 $\to$ 3.0), then passes the
    image and draft text through a Gemini 2.5 Flash multimodal polish pass.
  \item \textbf{Agent 3 --- Video Generator:} Produces per-scene video clips using
    Veo 3.1 Fast, with Chirp3-HD Text-to-Speech voiceover embedded for lip sync.
    Five scenes are generated in parallel via a \texttt{ThreadPoolExecutor} (three
    workers) and assembled into a 9:16 portrait MP4 using MoviePy. Chirp3-HD audio
    replaces Veo's native track to support SSML-controlled pacing.
  \item \textbf{Agent 4 --- Quality Control:} Scores output across seven criteria
    (accuracy, tone, engagement, compliance, platform fit, video script, and video
    visuals) using Gemini 2.5 Flash. Content scoring below 7 is flagged for revision;
    at or above 7 it is cleared for publication.
  \item \textbf{Agent 5 --- Scenario Resolver (Optional):} Provides deterministic
    financial guidance for specific F-1/J-1 student scenarios, operating independently
    from the main content pipeline.
\end{itemize}

The FastAPI backend exposes nine REST endpoints, including an asynchronous video job
system with browser-side polling via job ID. The React frontend stores pipeline state
in \texttt{localStorage} for non-volatile phase recovery across page reloads.

% ─────────────────────────────────────────────────────────────────────────────
\section{Features \& Functionality}
% ─────────────────────────────────────────────────────────────────────────────

Both sources agree on the core content generation capability. The slide deck emphasizes
the user-facing narrative; the repository confirms the full technical implementation.
Concrete capabilities include:

\begin{itemize}[leftmargin=1.5em]
  \item Topic selection across seven predefined civic and financial categories
  \item Platform toggle between LinkedIn and Instagram with persona-specific tone
    adaptation
  \item Live source retrieval with Google Search Grounding, returning cited results
  \item AI-generated post text, hashtags, and call-to-action copy
  \item Avatar editor: a customizable AI persona image generated via Imagen and
    persisted to the backend as a canonical avatar injected into content prompts
  \item Storyboard editor: reviewable and editable five-scene video brief before
    video generation is triggered
  \item Asynchronous video generation with browser-side polling and MP4 download
  \item Publishing queue with draft, scheduled, and published states, and optional
    webhook notifications on publish events
  \item Audience personalization sidebar with segment, tone, and platform style
    fields, persisted in \texttt{localStorage}
\end{itemize}

% ─────────────────────────────────────────────────────────────────────────────
\section{Known Limitations \& Risks}
% ─────────────────────────────────────────────────────────────────────────────

The slide deck explicitly identifies four risk categories, all of which find
technical correlates in the repository.

\begin{itemize}[leftmargin=1.5em]
  \item \textbf{Information accuracy:} Generated posts may contain wrong or outdated
    information. Agent 1 pulls live data but does not validate it against authoritative
    ground truth; Agent 4 scores accuracy using a language model judgment, which is
    itself fallible.
  \item \textbf{Overgeneralization:} The system uses a single persona configuration
    (\texttt{config/persona.py}) that may not capture the diversity of international
    student backgrounds, visa types, or countries of origin.
  \item \textbf{Sensitive topic handling:} Content on housing, immigration, and safety
    carries implicit advice risk. Explicit content filters beyond Agent 4's compliance
    criterion are not yet implemented.
  \item \textbf{Generic content:} Polished formatting does not guarantee relevance.
    The in-memory vector store does not persist or learn from user feedback between
    sessions.
  \item \textbf{Publishing stubs:} LinkedIn and Instagram adapters are currently
    placeholder implementations; no live publishing is active.
    \emph{[NOT IN SOURCES: production publishing integration timeline]}
\end{itemize}

% ─────────────────────────────────────────────────────────────────────────────
\section{Roadmap \& Next Steps}
% ─────────────────────────────────────────────────────────────────────────────

The slide deck defines a Trust-Ready MVP roadmap. The repository suggests additional
near-term technical work that complements it.

\textbf{From the slide deck:}
\begin{itemize}[leftmargin=1.5em]
  \item Embed source citations visibly in every published post draft (Agent 1 already
    returns cited sources; UI-level surfacing is
    \emph{[NOT IN SOURCES: current implementation status]})
  \item Human approval gate before any content is published
  \item User testing with real international students
  \item Improved content filters for sensitive topics (housing, immigration, safety)
  \item Platform expansion to TikTok, WhatsApp, and newsletters
\end{itemize}

\textbf{Additional work implied by the repository:}
\begin{itemize}[leftmargin=1.5em]
  \item Activate live LinkedIn and Instagram publishing adapters
  \item Implement persistent vector storage for cross-session retrieval quality
  \item Expand beyond the single \texttt{PERSONA\_CONFIG} to support multiple student
    segments and visa categories
\end{itemize}

% ─────────────────────────────────────────────────────────────────────────────
\section{Team}
% ─────────────────────────────────────────────────────────────────────────────

\begin{itemize}[leftmargin=1.5em]
  \item Joaquin Sardon
  \item Kulbir Kaur
  \item Melissa Lee
  \item Shivani Kabra
  \item Gerald Velasquez
\end{itemize}

\emph{[NOT IN SOURCES: individual role assignments or GitHub contributor breakdown]}

% ─────────────────────────────────────────────────────────────────────────────
\section{Conclusion}
% ─────────────────────────────────────────────────────────────────────────────

The International Student AI Influencer demonstrates a credible end-to-end AI content
pipeline that addresses a genuine information access gap. The system's architecture---
five specialized agents, multi-modal output across text, image, and video, and a modern
React and FastAPI interface---is technically ambitious for a hackathon context and
establishes a solid foundation for a trust-ready production tool. The team has correctly
identified the central challenges of accuracy, generalization, and sensitive content
handling, and the roadmap reflects a thoughtful path toward a responsible, student-tested
deployment.

\end{document}
